
\section{Related Work}
\vspace{2pt}\noindent\textbf{Cellular Network Security.}
A large body of research has examined the security of cellular networks across different protocol layers and interfaces. At the UE and RAN side, prior work has investigated baseband vulnerabilities~\cite{kim2021basespec, tu2024logic, maier2020basesafe}, lower-layer attacks~\cite{rupprecht2019breaking, yang2019hiding}, eavesdropping~\cite{rupprecht2020call}, privacy violations~\cite{hussain2019privacy}, and protocol compliance testing~\cite{kim2019touching, hussain2021noncompliance, park2022doltest, khandker2024astra, yang2024oranalyst}.
However, these efforts predominantly target the air interface between UE and RAN or the NAS/RRC signaling between UE and core, leaving the \textit{internal interfaces} of the core network relatively underexplored.
Recent work has started to examine core-network security from specific angles:
formal verification of access-control mechanisms~\cite{akon2023formal, akon2025control},
provenance-based attack detection~\cite{pacherkar2024prov5gc},
context-integrity testing of the LTE core~\cite{son2025citesting},
and GTP-U tunnel abuse that demonstrates external attackers can reach internal CN interfaces~\cite{zhang2025invade, shaik2025uncovering}.
Motivated by the growing reachability of internal interfaces in cloud-native deployments, our work takes a complementary direction: we systematically study implementation vulnerabilities across PFCP and GTP-C protocols and characterize \textit{iTrue} as a recurring type of security flaw, where core-network components blindly trust internal messages and skip critical validations, leading to severe consequences including DoS and session hijacking.

\vspace{2pt}\noindent\textbf{Automated Vulnerability Discovery in Cellular Networks.}
Fuzzing has been applied to cellular implementations from UE-RAN side testing~\cite{kim2019touching, park2022doltest, garbelini2022towards} to core-network fuzzing~\cite{bennett2024ransacked,dong2025corecrisis, sun20255gc, yang2024feedback}, but these approaches predominantly uncover memory-safety bugs and crash-inducing issues, while many high-impact core-network failures are logic vulnerabilities that do not manifest as crashes.
Specification-based analysis mines vulnerabilities from 3GPP documents through NLP~\cite{chen2021bookworm}, change-request analysis~\cite{chen2022seeing}, automated reasoning~\cite{chen2023sherlock}, undefined-behavior studies~\cite{klischies2023instructions}, and protocol state-machine synthesis~\cite{alishtiaq2024hermes, rahman2024cellularlint, xie2025cellsecinspector}; 
however, these methods target design flaws at the specification level and cannot automatically validate whether identified issues are exploitable in real implementations.
In contrast, iFinder distills recurring vulnerability patterns from previously discovered core network flaws and uses them to guide an LLM-driven multi-agent system, enabling efficient discovery of both memory-safety and deep logic vulnerabilities while closing the loop with code-specification cross-checking and runtime-verified PoC generation.

\vspace{2pt}\noindent\textbf{Multi-Agent for Vulnerability Discovery.}
LLMs have recently been applied to automate vulnerability discovery through agentic and multi-agent pipelines. 
Previous works have used LLMs to enhance existing static-analysis tools, e.g., improving the efficiency of CodeQL~\cite{wang2026cqllm} or customizing CodeQL for JavaScript vulnerability detection~\cite{ghebremichael2026multi}.
Others build general-purpose multi-agent pipelines to discover bugs in general programs~\cite{wu2026mulvul, liu2026synthesizing}.
However, these systems target general-purpose code and are not tailored to stateful protocols, so they are ineffective for 5G core networks and only report vulnerability candidates without confirming exploitability.
In contrast, \textit{iFinder} employs a pattern-guided multi-agent system that performs code-specification cross-checking to recover procedure context and iteratively refines PoCs with feedback, enabling it to detect iTrues in 5G core network implementations and confirm their exploitability.


\vspace{2pt}\noindent\textbf{Input Validation and Specification-Driven Testing.}
Some iTrues are related to input-validation vulnerabilities, a broad class that includes XSS~\cite{wassermann2008static}, SQL injection~\cite{al2023sqirl}, and, more recently, prompt injection~\cite{greshake2023not}, each with distinct root causes and mitigation strategies.
However, many other iTrues do not fall under input validation: they are triggered by well-formed inputs that arrive in incorrect system states. Because PFCP and GTP-C are long-lived, stateful protocols, inconsistencies between internal states and otherwise valid inputs can introduce security risks.
For example, a duplicate PDR ID may be syntactically valid yet still cause session hijacking. This distinction differentiates many iTrues from input-validation vulnerabilities.
Moreover, although iTrues are not fundamentally new, their detection cannot be effectively addressed by existing techniques. Specification-driven approaches, including model-based testing~\cite{utting2012taxonomy}, primarily target design flaws, while stateful protocol testing~\cite{ba2022stateful} and stateful network fuzzing~\cite{pham2020aflnet} often rely on crash-based signals and therefore struggle to capture logic flaws that do not manifest as crashes.













